% 原书物理页=印刷页13--16；1.5 全节（1.5.8 及其证明位于第 16 页顶部）.
\section{Lie 群的 Lie 代数}\label{sec:lie-1-5}

我们在此关心 Lie 群的微分性质. 特别地，我们将给 Lie 群在其单位元处的切空间 $T_eG$ 赋予一个 $\mathbb{R}$-Lie 代数结构.

\begin{numpara}[合成运算与求逆映射的微分]\label{np:lie-1-5-1}
设 $G$ 是一个 Lie 群，其合成运算记作 $m$，``求逆映射''记作 $i$，单位元为 $e$. 把合成运算 $m$ 在 $(e,e)$ 处微分，便得到一个 $\mathbb{R}$-线性映射 $d_{(e,e)}m:T_eG\times T_eG\longrightarrow T_eG$. $m$ 的结合律立即蕴含 $d_{(e,e)}$ 是一个满足结合律的运算. 对关系式 $m(g,e)=g$ 求微分可知，$0$ 是这一运算的单位元. 最后，把 $i$ 在点 $e$ 处微分，便得到 $T_eG$ 的一个自同态 $d_ei$，我们立即看出它给出运算 $T_eG$ 的``求逆映射''，因此这一运算是一个群运算.
这一运算是什么呢？正是 $T_eG$ 这个 $\mathbb{R}$-向量空间中的加法. 事实上
\[
d_{(e,e)}m(X,Y)=d_e[x\mapsto m(x,e)](X)+d_e[y\mapsto m(e,y)](Y)=d_e\operatorname{id}_G(X)+d_e\operatorname{id}_G(Y)=X+Y.
\]
因此，对任意 $X\in T_eG$ 都有 $d_ei(X)=-X$.
\end{numpara}

\begin{numpara}[一个 Lie 群的伴随表示]\label{np:lie-1-5-2}
对于元素 $x\in G$，记
\[
\operatorname{Int}_x:\begin{array}{r@{\ \to\ }l} G & G\\ y & xyx^{-1}\end{array}\quad\text{且}\quad \operatorname{Ad}_x:=d_e(\operatorname{Int}_x):T_eG\longrightarrow T_eG.
\]
由于 $\operatorname{Int}_x\circ\operatorname{Int}_{x'}=\operatorname{Int}_{xx'}$，我们同样有 $\operatorname{Ad}_x\circ\operatorname{Ad}_{x'}=\operatorname{Ad}_{xx'}$，从而 $\operatorname{Ad}_x$ 可逆，并且映射
\[
\begin{aligned}
\operatorname{Ad}:\quad &G\ \longrightarrow\ \mathrm{GL}(T_eG)\\
&x\ \longmapsto\ \operatorname{Ad}_x
\end{aligned}
\]
是一个群同态. 由于映射 $G\times G\longrightarrow G,\ (x,y)\mapsto xyx^{-1}$ 是光滑的，对其微分得到的映射 $G\times T_eG\longrightarrow T_eG,\ (x,Y)\mapsto\operatorname{Ad}_x(Y)$ 也是光滑的，因而态射 $\operatorname{Ad}$ 是光滑的. 我们把它称为 $G$ 的伴随表示. 它的底空间因此就是 $e$ 处的切空间.
\end{numpara}

\begin{numpara}[括号的构造]\label{np:lie-1-5-3}
由于 $\mathrm{GL}(T_eG)$ 是 $\mathbb{R}$-向量空间 $\operatorname{End}_{\mathbb{R}}(T_eG)$ 的一个开集，我们可以把后者等同于 $\mathrm{GL}(T_eG)$ 在任一点处的切空间，特别是在点 $e$ 处. 把 $\operatorname{Ad}$ 在点 $e$ 处微分，便得到一个 $\mathbb{R}$-线性映射
\[
\begin{aligned}
\operatorname{ad}:=d_e\operatorname{Ad}:\quad &T_eG\ \longrightarrow\ \operatorname{End}_{\mathbb{R}}(T_eG)\\
&X\ \longmapsto\ \operatorname{ad}_X
\end{aligned}
\]
剩下只需用
\[
[X,Y]:=\operatorname{ad}_X(Y),\qquad \text{其中 } X,Y\in T_eG
\]
来定义 Lie 括号.
由构造可知，这确实是一个双线性映射 $L\times L\longrightarrow L$\footnote{译注：原文如此；此处 $L$ 应理解为 $T_eG$.}，但尚需证明 Lie 括号的公理 i) 与 ii).
\end{numpara}

\begin{proposition}\label{prop:lie-1-5-4}
上面在 $T_eG$ 上定义的 Lie 括号是反对称的.
\end{proposition}

\begin{proof}
考虑\newterm{lie-commutator}{换位子}{commutator}映射 $c:G\times G\longrightarrow G,\ (x,y)\mapsto xyx^{-1}y^{-1}$. 这是一个光滑映射，它把 $G\times\{e\}$ 与 $\{e\}\times G$ 都送到 $\{e\}$.
将 $c$ 在点 $e$ 处对 $y$ 求微分，得到光滑映射 $G\times T_eG\longrightarrow T_eG,\ (x,Y)\mapsto\operatorname{Ad}_x(Y)-Y$.（注意 $c(x,y)=\operatorname{Int}_x(y)i(y)$.）再将这一映射在点 $e$ 处对 $x$ 求微分，便得到双线性映射
\[
\partial_x\partial_yc:T_eG\times T_eG\longrightarrow T_eG,\quad (X,Y)\mapsto\operatorname{ad}_X(Y).
\]
现在，先在点 $e$ 处对 $x$ 求 $c$ 的微分，得到光滑映射 $T_eG\times G\longrightarrow T_eG,\ (X,y)\mapsto X-\operatorname{Ad}_y(X)$.（注意 $c(x,y)=x\operatorname{Int}_y(i(x))$.）再将这一映射在点 $e$ 处对 $y$ 求微分，便得到双线性映射
\[
\partial_y\partial_xc:T_eG\times T_eG\longrightarrow T_eG,\quad (X,Y)\mapsto-\operatorname{ad}_Y(X).
\]
由于我们的映射都是光滑的，关于偏导数可交换性的 Schwarz 定理保证 $\partial_x\partial_yc=\partial_y\partial_xc$，由此便得到所需的反对称性.
\end{proof}

在证明 Jacobi 恒等式之前，我们先来验证括号构造的函子性.

\begin{proposition}\label{prop:lie-1-5-5}
设 $\varphi:G\longrightarrow G'$ 是一个 Lie 群态射.
\begin{enumerate}[label=\roman*)]
\item 对任意 $x\in G$，有 $d_e\varphi\circ\operatorname{Ad}_x=\operatorname{Ad}_{\varphi(x)}\circ d_e\varphi$.
\item 对任意 $X\in T_eG$，有 $d_e\varphi\circ\operatorname{ad}_X=\operatorname{ad}_{d_e\varphi(X)}\circ d_e\varphi$.
\end{enumerate}
\end{proposition}

\begin{proof}
i) 由于 $\varphi$ 是群的态射，对任意 $x\in G$ 有等式 $\varphi\circ\operatorname{Int}_x=\operatorname{Int}_{\varphi(x)}\circ\varphi$. 在点 $e$ 处对这一等式求微分，即得 i).

ii) 等式 i) 是光滑函数 $G\longrightarrow\operatorname{Hom}_{\mathbb{R}}(T_eG,T_{e'}G')$ 之间的等式. 在点 $e$ 处对这一等式求微分，即得 ii).
\end{proof}

我们把这一命题应用于 $G=G'$ 且 $\varphi=\operatorname{Int}_g$（$g\in G$ 为某一元素）的情形. 此时 ii) 给出 $\operatorname{Ad}_g\circ\operatorname{ad}_X=\operatorname{ad}_{\operatorname{Ad}_g(X)}\circ\operatorname{Ad}_g$，这也可以写成
\[
\operatorname{Ad}_g([X,Y])=[\operatorname{Ad}_gX,\operatorname{Ad}_gY],\qquad \forall Y\in T_eG.
\]
\begin{corollary}\label{cor:lie-1-5-6}
上面在 $T_eG$ 上定义的 Lie 括号满足 Jacobi 恒等式.
\end{corollary}

\begin{proof}
在点 $g=e$ 处对上面的恒等式求微分，便对任意 $Z\in T_eG$ 得到
\begin{align*}
\operatorname{ad}_Z([X,Y])\ &=\ d_{(e,e)}\left((g,h)\mapsto[\operatorname{Ad}_gX,\operatorname{Ad}_hY]\right)(Z,Z)\\
&=\ [\operatorname{ad}_ZX,\operatorname{Ad}_eY]+[\operatorname{Ad}_eX,\operatorname{ad}_ZY]\\
&=\ [\operatorname{ad}_Z(X),Y]+[X,\operatorname{ad}_Z(Y)]
\end{align*}
这等价于 Jacobi 恒等式（至少等到我们证明了括号的反对称性之后就等价）.
\end{proof}

\begin{numpara}[$G$ 的 Lie 代数]\label{np:lie-1-5-7}
我们把所得到的结果总结如下：
\begin{enumerate}[label=\roman*)]
\item 配有括号~\ref{np:lie-1-5-3}~的切空间 $T_eG$ 是一个 $\mathbb{R}$-Lie 代数，我们把它记作 $\operatorname{Lie}(G)$.
\item 任何 Lie 群态射 $\varphi:G\longrightarrow G'$ 都通过在单位元处的微分，诱导出一个 Lie 代数态射 $d\varphi:\operatorname{Lie}(G)\longrightarrow\operatorname{Lie}(G)$\footnote{译注：原文如此；依上下文应为 $\operatorname{Lie}(G')$.}.
\item $\operatorname{Lie}(G\times G')=\operatorname{Lie}(G)\times\operatorname{Lie}(G')$（显然）.
\item $\operatorname{Lie}(G)=\operatorname{Lie}(G^0)$（显然）.
\item $\operatorname{Lie}(G)$ 是伴随表示 $\operatorname{Ad}_G:G\longrightarrow\operatorname{Lie}(G)$ 的空间.
\end{enumerate}
\end{numpara}

\noindent\textbf{例.}
\begin{enumerate}[label=\roman*)]
\item \newterm{lie-discrete-group}{离散群}{discrete group}的 Lie 代数是零代数.
\item $\operatorname{Lie}(\mathbb{R})=\operatorname{Lie}(\mathbb{R}^{\times})=\operatorname{Lie}(S^1)=\mathbb{R}$（配备零 Lie 括号）. 更一般地，\newterm{lie-abelian}{阿贝尔}{abelian} Lie 群的 Lie 括号为零.
\item $\operatorname{Lie}(\mathrm{GL}_n(\mathbb{R}))=\mathfrak{gl}_n(\mathbb{R})$（记号见上文），并且伴随表示由共轭给出. 事实上，$I_n$ 处的切空间等同于 $\mathrm{M}_n(\mathbb{R})$. 设 $x,y\in\mathrm{GL}_n(\mathbb{R})$，并设 $y=I_n+Y$，其中 $\|Y\|<1$. 等式 $xyx^{-1}=I_n+xYx^{-1}$ 表明 $\operatorname{Ad}_x$ 就是 $x$ 在 $\mathrm{M}_n(\mathbb{R})$ 上取共轭. 再设 $x=I_n+X$，其中 $\|X\|<1$. 于是有 $\operatorname{Ad}_x(Y)=(I_n+X)Y(I_n-X)+o(X)=Y+[X,Y]+o(X)$，这表明 $\operatorname{ad}_X(Y)=[X,Y]$.
\item 更内蕴地，$\operatorname{Lie}(\mathrm{GL}(V))=\mathfrak{gl}(V)$.
\end{enumerate}

\begin{exercise*}
证明行列式 $\det:\mathrm{GL}_n(\mathbb{R})\longrightarrow\mathbb{R}^{\times}$ 的微分是迹 $\operatorname{tr}:\mathfrak{gl}_n(\mathbb{R})\longrightarrow\mathbb{R}$.
\end{exercise*}

下面给出群的性质与其 Lie 代数的性质之间相互作用的一个例子.

\begin{proposition}\label{prop:lie-1-5-8}
设 $\varphi:G\longrightarrow G'$ 是一个 Lie 群态射.
\begin{enumerate}[label=\roman*)]
\item $d\varphi$ 是满射 $\Leftrightarrow$ $\varphi(G)\supset G'^{0}$，
\item $d\varphi$ 是单射 $\Leftrightarrow$ $\operatorname{Ker}(\varphi)$ 是离散的，
\item 若 $G'$ 连通，则 $d\varphi$ 是双射 $\Leftrightarrow$ $\varphi$ 是一个覆盖.
\end{enumerate}
\end{proposition}

\begin{proof}
我们已经看到，i) 可由常秩定理得出. ii) 也可由它得出：若 $d_{e}\varphi$ 是单射，则它在 $e$ 的某个邻域内是单射，而常秩定理给出 $e$ 的一个邻域，使得 $\varphi$ 在其上的限制是单射. 于是核 $\operatorname{Ker}(d\varphi)$ 是离散的. 逆命题也是显然的. 最后，iii) 可由 i) 与 ii) 得出：设 $\mathcal{V}$ 是 $e$ 的一个邻域，满足 $\mathcal{V}\cap\operatorname{Ker}(\varphi)=\{e\}$. 由乘法运算与求逆映射的连续性，可以找到 $e$ 的一个邻域 $\mathcal{U}$，使得 $\mathcal{U}.\mathcal{U}^{-1}\subset\mathcal{V}$. 于是对 $x\neq y\in\operatorname{Ker}(\varphi)$，有 $x\mathcal{U}\cap y\mathcal{U}=\varnothing$. 令 $\mathcal{U}':=\varphi(\mathcal{U})$. 它是 $G'$ 中 $e'$ 的一个开邻域（因为 $d\varphi$ 是双射），并且 $\varphi^{-1}(\mathcal{U}')\simeq\operatorname{Ker}(\varphi)\times\mathcal{U}$. 经平移便得到：对一切 $g'\in G'$ 与 $g\in\varphi^{-1}(g')$ 都有 $\varphi^{-1}(g'\mathcal{U}')\simeq\operatorname{Ker}(\varphi)\times g\mathcal{U}$，这表明 $\varphi$ 是一个覆盖.
\end{proof}
